% AUTHOR: Amlal El Mahrouss
% PURPOSE: D01: The T-Function.


\documentclass[11pt, a4paper]{article}
\usepackage{graphicx}
\usepackage{listings}
\usepackage{amsmath,amssymb,amsthm}
\usepackage{xcolor}
\usepackage{hyperref}
\usepackage[margin=0.5in,top=1in,bottom=1in]{geometry}

\title{The T-Function.}
\author{Amlal El Mahrouss\\\texttt{amlal@nekernel.org}}
\date{February 2026}

\begin{document}

\bf \maketitle

\begin{center}
	\rule[0.01cm]{17cm}{0.01cm}
\end{center}

\section{Abstract}
This paper covers a function $\Theta(\Delta t)$, which we will use to solve a set of lemmas in $\mathbb{C}$.

\section{Definitions}

\subsection{Definition of $\Theta(\Delta t)$}
Let the T-Function:
$\Theta(\Delta t, \Delta\phi)=\frac{\Delta\phi}{\Delta t}$
such that $\Theta(\Delta t) \in \mathbb{C}$ if and only if $\Theta(\Delta t, \Delta\phi) > \frac{\Delta\phi}{\Delta t}$

\subsection{Lemma 1}

Proof for $\Theta(\Delta t, \Delta\phi) > \frac{\Delta\phi}{\Delta t}$ when $\forall \Delta t = 4i, \quad \Delta\phi = \ln(e \times i)$ then
\begin{equation}
    \Theta(\Delta t, \Delta\phi) > \frac{\Delta\phi}{{\Delta t}}
\end{equation}
Hence
\begin{equation}
(\Theta(\Delta t, \Delta\phi) \times \Delta t) < \Delta\phi
\end{equation}
Then
\begin{equation}
\Theta(\Delta t, \Delta\phi) \times \Delta t \in \mathbb{C}
\end{equation}
Then
\begin{equation}
y = \ln(e \times i) \to \ln (e(i \times t))
\end{equation}
Such that we then multiply $\Delta t$ from both sides:
\begin{equation}
    y = |\Theta(\Delta t, \Delta\phi) \times \Delta (n \times i)| > \ln (e \times (i \times t))
\end{equation}
Hence

\begin{equation}
    \forall y \in \mathbb{C}, \quad y > \ln(e \times i)
\end{equation}
This completes the proof.

\end{document}
